% Unit 107 terjemahan Bahasa Indonesia dari chapter8.tex baris 175--281.
% Sumber: Methods of Algebra, Volume 2, commit
% 9a5803ff77dd3257484cb177f851a73770a59dd3 (CC BY 4.0).
% stable-unit-id: o014.aljabr2.chapter8.simplicial-sets
% Cakupan: Bagian 8.2 lengkap, ``Himpunan Simplisial'';
% berhenti tepat sebelum Bagian 8.3 pada baris 282.

% segment-id: o014.li2.chapter8.simplicial-sets.h001
\section{Himpunan Simplisial}\label{sec:simplicial-set}
\hypertarget{o014.aljabr2.chapter8.simplicial-sets}{ }
% segment-id: o014.li2.chapter8.simplicial-sets.p001
Objek simplisial dari Definisi~\ref{def:simplicial-obj} dalam kasus
$\mathcal{C}=\cate{Set}$ disebut himpunan simplisial. Dengan semesta
Grothendieck yang telah dipilih, arti tepat $\cate{Set}$ dalam kerangka
buku ini adalah kategori semua himpunan kecil; rincian tersebut tidak
memengaruhi isi bagian ini.

% segment-id: o014.li2.chapter8.simplicial-sets.d001
\begin{definition}
	\index{himpunan simplisial@himpunan simplisial (simplicial set)}
	\index{himpunan kosimplisial@himpunan kosimplisial (cosimplicial set)}
	Objek simplisial dalam kategori himpunan $\cate{Set}$ disebut
	\emph{himpunan simplisial}; objek kosimplisial di dalamnya disebut
	\emph{himpunan kosimplisial}.
\end{definition}

% segment-id: o014.li2.chapter8.simplicial-sets.p002
Menurut Definisi~\ref{def:simplicial-obj}, semua himpunan simplisial (atau
kosimplisial) membentuk kategori
$\cate{sSet}:=\cate{Set}^{\simpDelta^{\opp}}$ (atau
$\cate{csSet}:=\cate{Set}^{\simpDelta}$).
\index[sym1]{sSet@$\cate{sSet}, \cate{csSet}$}

% segment-id: o014.li2.chapter8.simplicial-sets.d002
\begin{definition}\label{def:standard-simplicial-set}
	\index[sym1]{Delta-n@$\Delta^n$}
	Untuk setiap $n\in\Z_{\geq0}$, tuliskan $\Delta^n$ untuk himpunan
	simplisial yang ditentukan oleh
	$\Hom_{\simpDelta}(\mathord\cdot,[n]):
	\simpDelta^{\opp}\to\cate{Set}$. Objek ini disebut \emph{simpleks-$n$
	standar}.
\end{definition}

% segment-id: o014.li2.chapter8.simplicial-sets.p003
Jadi, $(\Delta^n)_m=\Hom_{\simpDelta}([m],[n])$ adalah himpunan semua
pemetaan monoton $[m]\to[n]$. Pemetaan
$(\Delta^n)_{m'}\to(\Delta^n)_m$ yang diinduksi oleh pemetaan monoton
$\phi:[m]\to[m']$ tepat merupakan penarikan balik pemetaan
$\phi^*:f\mapsto f\phi$.

% segment-id: o014.li2.chapter8.simplicial-sets.eg001
\begin{example}\label{eg:boundary-horn}
	\index[sym1]{partial-Delta-n@$\partial \Delta^n$}
	\index[sym1]{Lambda-n-k@$\Lambda^n_k$}
	\index{tanduk@tanduk (horn)}
	Kita perkenalkan dua subobjek umum dari simpleks-$n$ standar $\Delta^n$.
	\begin{description}
		\item[Batas] Definisikan subfunktor $\partial\Delta^n$ dari $\Delta^n$
		melalui
		% segment-id: o014.li2.chapter8.simplicial-sets.q001
		\[ (\partial \Delta^n)_m := \left\{ f: [m] \to [n] \; \text{monoton}, \Image(f) \neq [n] \right\}, \quad m \in \Z_{\geq 0}. \]
		Bila $m<n$, berlaku
		$(\partial\Delta^n)_m=(\Delta^n)_m$; bila $m\geq n$, elemen
		$(\partial\Delta^n)_m$ diperoleh melalui degenerasi berulang dari
		$(\Delta^n)_h$ untuk $0\leq h<n$.
		\item[Tanduk] Untuk $0\leq k\leq n$, definisikan subfunktor
		$\Lambda^n_k$ dari $\Delta^n$ melalui
		% segment-id: o014.li2.chapter8.simplicial-sets.q002
		\[ (\Lambda^n_k)_m := \left\{ f: [m] \to [n] \;\text{monoton}, \Image(f) \not\supset [n] \smallsetminus \{k\} \right\}, \quad m \in \Z_{\geq 0}. \]
		Kita sebut $k$ simpul tanduk tersebut.
	\end{description}

	Setiap penarikan balik $\phi^*:f\mapsto f\phi$ jelas mempertahankan
	kondisi di atas, sehingga keduanya merupakan subfunktor. Latihan bab ini
	akan menunjukkan bagaimana $\Lambda^n_i$ (atau $\partial\Delta^n$) dapat
	dinyatakan secara tepat sebagai penempelan $n$ (atau $n+1$) salinan
	$\Delta^{n-1}$.
\end{example}

% segment-id: o014.li2.chapter8.simplicial-sets.p004
Himpunan simplisial konstan yang diberikan oleh himpunan kosong,
$\mathrm{const}(\emptyset)$, disingkat $\emptyset$. Dengan demikian,
$\partial\Delta^0=\emptyset$.

% segment-id: o014.li2.chapter8.simplicial-sets.p005
Dalam \S\ref{sec:geom-realization}, kita akan memberikan interpretasi intuitif
himpunan simplisial. Khususnya, akan digambarkan citra geometris ketiga objek
$\Delta^n\supset\Lambda^n_i\supset\partial\Delta^n$.

% segment-id: o014.li2.chapter8.simplicial-sets.p006
Lema Yoneda (Teorema~\ref{prop:Yoneda}) langsung memberikan isomorfisme
kanonik berikut, dengan $n\in\Z_{\geq0}$ dan
$X\in\Obj(\cate{sSet})$:
% segment-id: o014.li2.chapter8.simplicial-sets.q003
\begin{equation}\begin{gathered}
	\begin{tikzcd}[row sep=tiny]
		\Hom_{\cate{sSet}}\left( \Delta^n, X \right) \arrow[r, "\sim"] & X_n \\
		\phi \arrow[mapsto, r] & \phi([n])(\identity_{[n]})
	\end{tikzcd} \\
	\phi([n]): \End([n]) = \Delta^n([n]) \to X([n]) =: X_n.
\end{gathered}\end{equation}

% segment-id: o014.li2.chapter8.simplicial-sets.d003
\begin{definition}\label{def:simplex}
	\index{simpleks@simpleks (simplex)}
	\index{simpleks!tak terdegenerasi@tak terdegenerasi (non-degenerate)}
	Elemen $X_n$ disebut \emph{simpleks-$n$} dalam himpunan simplisial $X$.
	Jika $x\in X_n$ terletak dalam citra suatu $s_j$, ia disebut
	\emph{degenerat}; jika tidak, ia disebut \emph{nondegenerat}.
\end{definition}

% segment-id: o014.li2.chapter8.simplicial-sets.p007
Sebagai latihan sederhana, periksa bahwa simpleks-$m$
nondegenerat dari $\Delta^n$ bersesuaian dengan injeksi monoton
$[m]\hookrightarrow[n]$; jumlahnya $\binom{n+1}{m+1}$.

% segment-id: o014.li2.chapter8.simplicial-sets.p008
Lebih lanjut, Lema Yoneda memastikan bahwa
$h_{\simpDelta}:[n]\mapsto\Delta^n$ mendefinisikan funktor yang penuh dan setia
$\simpDelta\to\cate{sSet}$. Setiap morfisme $f:[m]\to[n]$ dalam
$\simpDelta$ menginduksi $f:\Delta^m\to\Delta^n$. Memberikan subhimpunan
beranggota $m+1$, $\{k_0,\ldots,k_m\}\subset[n]$, sama artinya dengan
memberikan injeksi monoton $[m]\hookrightarrow[n]$. Morfisme yang
bersesuaian antara simpleks standar juga ditulis
% segment-id: o014.li2.chapter8.simplicial-sets.q004
\[ \Delta^m \xrightarrow{\{k_0, \ldots, k_m\}} \Delta^n . \]

% segment-id: o014.li2.chapter8.simplicial-sets.p009
Makna intuitif himpunan simplisial akan dijelaskan melalui realisasi
geometris dalam \S\ref{sec:geom-realization}. Sebelum itu, kita perkenalkan
serangkaian konstruksi abstrak, dimulai dengan join serta kerucut kiri dan
kanan himpunan simplisial.

% segment-id: o014.li2.chapter8.simplicial-sets.p010
Tuliskan $\cate{FinLin}_+$ untuk subkategori kategori himpunan terurut
total hingga (atau himpunan terurut linear hingga) $\cate{FinLin}$ yang
diperoleh dengan membuang himpunan kosong; morfismenya adalah pemetaan
monoton. Karena $\simpDelta$ merupakan kerangka dari $\cate{FinLin}_+$,
himpunan simplisial secara ekuivalen dapat dipandang sebagai funktor
$\cate{FinLin}_+^{\opp}\to\cate{Set}$.

% segment-id: o014.li2.chapter8.simplicial-sets.cv001
\begin{convention}
	Misalkan $J$ objek $\cate{FinLin}_+$. Setiap subhimpunan $I\subset J$
	secara otomatis merupakan himpunan terurut total hingga. Jika juga berlaku
	% segment-id: o014.li2.chapter8.simplicial-sets.q005
	\[ \forall (i, j) \in I \times J, \quad j \leq i \implies j \in I, \]
	$I$ disebut \emph{ruas awal} dari $J$ dan ditulis $I\sqsubset J$.
\end{convention}

% segment-id: o014.li2.chapter8.simplicial-sets.d004
\begin{definition}\label{def:simplicial-join}
	\index{himpunan simplisial!join@join}
	\emph{Join} $X\star X'$ dari himpunan simplisial $X$ dan $X'$
	didefinisikan sebagai funktor $\cate{FinLin}_+^{\opp}\to\cate{Set}$
	berikut:
	% segment-id: o014.li2.chapter8.simplicial-sets.q006
	\[ (X \star X')(J) = \bigsqcup_{I \sqsubset J} X(I) \times X'(J \smallsetminus I). \]
	Di sini (dan hanya di sini), $X(\emptyset)$ dan $X'(\emptyset)$
	didefinisikan sebagai himpunan satu titik. Untuk setiap pemetaan monoton
	$f:J_1\to J$ dan $I\sqsubset J$, tetapkan $I_1:=f^{-1}(I)$. Maka
	$I_1\sqsubset J_1$ dan terdapat pemetaan
	% segment-id: o014.li2.chapter8.simplicial-sets.q007
	\[ X(I) \times X'(J \smallsetminus I) \to X(I_1) \times X'(J_1 \smallsetminus I_1). \]
	Dengan demikian diperoleh morfisme terinduksi
	$f^*:(X\star X')(J)\to(X\star X')(J_1)$, yang menentukan struktur
	simplisial pada $X\star X'$.
\end{definition}

% segment-id: o014.li2.chapter8.simplicial-sets.p011
Terdapat morfisme kanonik $X\to X\star Y\leftarrow Y$. Definisi tersebut
juga dapat ditulis dalam notasi sebelumnya sebagai
% segment-id: o014.li2.chapter8.simplicial-sets.q008
\[ (X \star Y)_n = X_n \sqcup Y_n \sqcup \bigsqcup_{j+k = n-1} X_j \times Y_k . \]
Morfisme $d_i$ dan $s_i$ diperoleh langsung dari morfisme yang bersesuaian
pada $X$ dan $Y$. Sebagai contoh, untuk $d_i:(X\star Y)_n\to(X\star Y)_{n-1}$,
pembatasannya pada subhimpunan $X_n$ dan $Y_n$ adalah $d_i$ semula,
sedangkan pada subhimpunan $X_j\times Y_k$ diberikan oleh
% segment-id: o014.li2.chapter8.simplicial-sets.q009
\begin{equation}\label{eqn:join-formulas}\begin{gathered}
		d_i(x, y) = \begin{cases}
			(d_i x, y), & i \leq j, \; j \neq 0 \\
			(x, d_{i-j-1} y), & i > j, \; k \neq 0,
		\end{cases} \\
		j=0 \implies d_0(x, y) = y \in Y_{n-1}, \\
		k=0 \implies d_n(x, y) = x \in X_{n-1}.
\end{gathered}\end{equation}

% segment-id: o014.li2.chapter8.simplicial-sets.pr001
\begin{proposition}\label{prop:join-nd}
	\index[sym1]{Xnd@$X^{\mathrm{nd}}$}
	Tuliskan $X_n^{\mathrm{nd}}\subset X_n$ untuk subhimpunan simpleks-$n$
	nondegenerat dari himpunan simplisial $X$
	(Definisi~\ref{def:simplex}). Untuk sebarang $X$ dan $Y$, berlaku
	% segment-id: o014.li2.chapter8.simplicial-sets.q010
	\[ (X \star Y)^{\mathrm{nd}}_n = X^{\mathrm{nd}}_n \sqcup Y^{\mathrm{nd}}_n \sqcup \bigsqcup_{j+k=n-1} X^{\mathrm{nd}}_j \times Y^{\mathrm{nd}}_k. \]
\end{proposition}
% segment-id: o014.li2.chapter8.simplicial-sets.pf001
\begin{proof}
	Misalkan $Z$ himpunan simplisial. Untuk sebarang himpunan terurut total
	hingga tak kosong $J$, elemen $x\in Z(J)$ degenerat jika dan hanya jika
	terdapat pemetaan monoton surjektif tetapi tidak injektif
	$f:J\twoheadrightarrow J_0$ sedemikian sehingga
	$x\in\Image[f^*:Z(J_0)\to Z(J)]$. Pernyataan selebihnya langsung mengikuti
	pengamatan ini.
\end{proof}

% segment-id: o014.li2.chapter8.simplicial-sets.d005
\begin{definition}\label{def:cone-simplicial}
	\index{himpunan simplisial!kerucut@kerucut (cone)}
	\emph{Kerucut kiri} dengan alas himpunan simplisial $X$ adalah
	$X^{\lhd}:=\Delta^0\star X$, sedangkan \emph{kerucut kanan} adalah
	$X^{\rhd}:=X\star\Delta^0$.
\end{definition}

% segment-id: o014.li2.chapter8.simplicial-sets.p012
Interpretasi intuitif kerucut diberikan dalam
Contoh~\ref{eg:cone-realization}. Latihan akan membahas lebih banyak sifat
operasi join.
